\chapter{The Frontend and the Single Host}
\label{ch:frontend}

Milestone~7 puts the last deployable --- the React frontend --- into
the cluster and claims the \code{/} path on \code{ts.local} that
Chapter~18's Ingress deliberately left free. It is a short milestone
in lines of YAML and a long one in concepts: it is the only image the
pipeline builds from a Dockerfile, the only pod that is not a JVM, and
the change that turns ``services reachable with \code{curl}'' into
``an application usable in a browser''.

\section{What a frontend is, to an operator}

A built React app is not a server. \code{npm run build} (TypeScript
check, then Vite bundling) produces a directory of static files:
one \code{index.html}, one content-hashed JavaScript bundle
(\code{index-CqZING8N.js}), one CSS file, a favicon. Nothing in it
runs on the server; every line of it runs in the visitor's browser,
which then calls \code{/api/...} over HTTP like any other client. So
``deploying the frontend'' means deploying a \emph{file server} with
those files inside --- nginx here --- and telling the edge router which
requests are files and which are API.

That framing decides everything else in this chapter. No config
server (the app has no properties to import), no Eureka (nobody calls
it by name), no secret (the browser can read anything the frontend
holds, so nothing secret may be there), and a resource footprint
smaller than the smallest JVM's class-loading overhead.

\section{Two stages, one Dockerfile}

Jib (Chapter~13) builds JVM images from \code{target/classes}; it
cannot run \code{npm}. The frontend therefore gets the classic
multi-stage Dockerfile --- and the reason it has two stages is the same
reason the JVM images use a JRE base rather than a JDK:

\begin{lstlisting}[language=dockerfile]
FROM node:24-alpine AS build            (*@\co{1}@*)
WORKDIR /app
COPY package.json package-lock.json ./  (*@\co{2}@*)
RUN npm ci --no-audit --no-fund
COPY . .
RUN npm run build

FROM nginx:1.27-alpine                  (*@\co{3}@*)
COPY nginx.conf /etc/nginx/conf.d/default.conf
COPY --from=build /app/dist /usr/share/nginx/html   (*@\co{4}@*)
EXPOSE 80
\end{lstlisting}

\begin{description}
  \item[\co{1}] The build stage has Node, the TypeScript compiler, and
  after \code{npm ci} roughly 300\,MB of \code{node\_modules}. None of
  it is needed to \emph{serve} the result.
  \item[\co{2}] Lockfile first, sources second --- Chapter~12's layer
  cache rule. \code{npm ci} installs exactly the versions pinned in
  \code{package-lock.json} and refuses to modify it; a source edit
  invalidates only the layers from \code{COPY . .} down, so the
  dependency download is skipped on every rebuild that does not touch
  \code{package.json}.
  \item[\co{3}] A fresh base. The final image is nginx plus static
  files: about 15\,MB, no Node, no compiler, no package manager --- a
  smaller attack surface and a faster pull, for free.
  \item[\co{4}] \code{--from=build} is the multi-stage copy: only
  \code{dist/} crosses from stage one to stage two. Everything else
  in stage one is discarded when the build ends.
\end{description}

\begin{decision}[Why not add Node to the Jenkins image and use Jib?]
Jib can wrap any directory into an image, so a ``Node in Jenkins, Jib
for packaging'' route exists. It was rejected because it hides the
most common image-building technique behind the least common one.
Every team that ships a frontend has this Dockerfile; being able to
read and write it is a job-market skill in itself, and having one
non-Jib image in the pipeline also proves the pipeline is not secretly
Jib-shaped. The cost is a Docker CLI in the Jenkins image and a
mounted daemon socket --- Section~\ref{sec:dockersock}.
\end{decision}

\section{nginx: the fallback that makes a SPA work}

\begin{lstlisting}[language=bash]
server {
    listen 80;
    root  /usr/share/nginx/html;
    location / {
        try_files $uri $uri/ /index.html;          (*@\co{1}@*)
    }
    location /assets/ {
        add_header Cache-Control "public, max-age=31536000, immutable"; (*@\co{2}@*)
    }
    location = /healthz { return 200 "ok\n"; }      (*@\co{3}@*)
}
\end{lstlisting}

\begin{description}
  \item[\co{1}] The line that every single-page app needs and every
  first deployment forgets. React Router (\code{BrowserRouter} in
  \code{main.tsx}) owns URLs like \code{/courses/42} purely in the
  browser: clicking a link never asks the server. But a refresh, a
  bookmark or a pasted link \emph{does} ask the server for
  \code{/courses/42}, no such file exists, and a plain file server
  answers 404. \code{try\_files} tries the literal path, then a
  directory, then falls back to \code{index.html} --- the app boots
  and its router reads the URL. The OAuth callback
  (\code{/oauth/callback}) relies on this too: Google's redirect is a
  cold request for a route that exists only in JavaScript.
  \item[\co{2}] Vite writes a content hash into every asset filename,
  so a changed file always has a new name. That makes ``cache forever''
  safe for \code{/assets/} --- and \emph{unsafe} for
  \code{index.html}, which must be re-fetched to learn the new names
  after a deploy. Caching the HTML is how users end up with a stale
  app that requests bundles the server no longer has.
  \item[\co{3}] A probe target that costs nothing and never touches
  disk (Chapter~17's probes need something to call).
\end{description}

Note what the file does \emph{not} do: proxy \code{/api}. On a
laptop, \code{vite.config.ts} proxies \code{/api} to the gateway on
\code{localhost:8080} so the dev server and the API look like one
origin. In the cluster that job belongs to Traefik --- nginx never sees
an API request at all.

\section{Claiming \code{/}: one origin}

\begin{lstlisting}[language=yaml]
kind: Ingress
metadata: { name: frontend }
spec:
  rules:
    - host: ts.local
      http:
        paths:
          - path: /                     (*@\co{1}@*)
            pathType: Prefix
            backend: { service: { name: frontend, port: { number: 80 } } }
\end{lstlisting}

\begin{description}
  \item[\co{1}] A catch-all, in a \emph{second} Ingress object next to
  the gateway's. Traefik merges every Ingress rule for a host and routes
  by \emph{longest matching prefix}: \code{/api/courses} matches both
  \code{/} and \code{/api}, and \code{/api} wins. The gateway's three
  prefixes therefore keep working unchanged, and everything else ---
  \code{/}, \code{/login}, \code{/courses/42}, \code{/assets/...} ---
  falls through to nginx. Two objects, one merged routing table;
  Chapter~20's ``apply never deletes'' still applies to each.
\end{description}

The consequence is the reason this milestone exists at all:
\textbf{one origin}. The browser loads the page from
\code{http://ts.local} and calls \code{http://ts.local/api/...} ---
same scheme, host and port. Compare Compose, where the page came from
\code{localhost:3000} and the API from \code{localhost:8080}: a
different port is a different origin, so every API call was preceded
by a CORS preflight and the gateway's \code{CorsConfig} had to list
\code{localhost:3000} explicitly. In the cluster \code{CorsConfig}
is simply never consulted --- there is no cross-origin request to
police. Cookies, the OAuth redirect, and any future
\code{SameSite} policy all get simpler for the same reason.

\begin{figure}[h]
\centering
\begin{tikzpicture}[node distance=5mm and 8mm]
  \node[boxdark] (b) {Browser};
  \node[box, right=22mm of b] (t) {Traefik};
  \node[box, right=34mm of t, yshift=9mm] (fe) {frontend\\\scriptsize nginx};
  \node[box, right=34mm of t, yshift=-9mm] (gw) {api-gateway};
  \draw[flow] (b) -- node[lbl,above]{\scriptsize Host: ts.local} (t);
  \draw[flow] (t.east) -- node[lbl,above,sloped,pos=0.5]{\scriptsize /} (fe.west);
  \draw[flow] (t.east) -- node[lbl,below,sloped,pos=0.5]{\scriptsize /api, /oauth2, /login/oauth2} (gw.west);
\end{tikzpicture}
\caption{Longest prefix wins: one host, two backends.}
\end{figure}

One property follows the host change. \code{OAuth2LoginSuccessHandler}
finishes Google login by redirecting the browser to
\code{app.oauth2.success-redirect-uri}, which config-server's
\code{auth-service.yml} sets to \code{localhost:3000/oauth/callback}
--- correct on a laptop, wrong in the cluster. \code{auth-service.yaml}
overrides it with the env var
\code{APP\_OAUTH2\_SUCCESS\_REDIRECT\_URI=http://ts.local/oauth/callback}
--- Chapter~16's precedence rule (env var beats config server) doing
the same job it did for the datasource URL.

\section{The manifest}

The Deployment is Chapter~15's shape with almost everything removed:
no \code{envFrom}, no \code{SPRING\_*} variables, no
\code{startupProbe} (nginx is up in milliseconds), and
\code{requests: \{ memory: 32Mi, cpu: 10m \}} --- one-twelfth of the
smallest JVM. \code{imagePullPolicy: Always} with the
\code{:latest} tag as the manifest default, overridden to the commit
tag by the pipeline exactly as for the JVM services.

\section{The pipeline grows a stage}
\label{sec:dockersock}

\begin{lstlisting}
stage('Frontend image') {
  steps {
    dir('Project/TeacherSupporter') {
      sh '''
        docker build -t ${REGISTRY}/frontend:${TAG} \
                     -t ${REGISTRY}/frontend:latest frontend   (*@\co{1}@*)
        docker push ${REGISTRY}/frontend:${TAG}
        docker push ${REGISTRY}/frontend:latest
      '''
    }
  }
}
\end{lstlisting}

\begin{description}
  \item[\co{1}] The build context is \code{frontend/}, trimmed by a
  \code{.dockerignore} that excludes \code{node\_modules} and
  \code{dist}: the Windows laptop's native modules must not be sent to
  the daemon, and the image builds its own. \code{frontend} joins
  \code{SERVICES} (kubectl's list) but not \code{MODULES} (Maven's) ---
  the Deploy stage's \code{set image} / \code{rollout status} loops
  handle it without change.
\end{description}

For \code{docker build} to work inside the Jenkins container, two
things change in Chapter~22's setup. The image gains the Docker
\emph{CLI} (\code{docker-ce-cli} --- the client only), and
\code{run.sh} mounts the host's \code{/var/run/docker.sock} into the
container. The CLI talks over that socket to the \emph{host's} daemon;
there is no daemon inside Jenkins and no Docker-in-Docker. The
resulting image lands in the host's image store, and
\code{docker push localhost:5000/...} reaches the registry over the
host network the container already shares.

\begin{pitfall}[permission denied on /var/run/docker.sock]
Symptom: the CLI is installed, the socket is mounted, and every
\code{docker} command fails with a permission error. Cause: the
socket is owned by the host's \code{docker} group --- GID 984 on
\code{ts-server} --- and that number means nothing inside the image,
where the \code{jenkins} user belongs to a group of the same name with
a different GID (or none). Fix in \code{run.sh}:
\code{--group-add "\$(stat -c \%g /var/run/docker.sock)"} adds the
host's numeric GID to the container's user at run time, without baking
a machine-specific number into the Dockerfile.
\end{pitfall}

\begin{decision}[The socket is root]
Mounting the Docker socket grants root-equivalent control of the host:
whoever can run \code{docker run -v /:/host} owns the machine. For a
single-user homelab whose Jenkins is reachable only over the tailnet
this is an accepted trade, stated out loud. Production builds images
without a daemon at all --- Kaniko or BuildKit running as an
unprivileged pod, or Jib for the JVM --- and the Chapter~27 migration
of Jenkins into the cluster is where this compromise would be paid
back.
\end{decision}

\begin{hardware}
The frontend pod is the cheapest in the namespace (32\,Mi requested),
but the \emph{build} is not: \code{npm ci} plus \code{tsc} plus Vite
runs near a CPU-minute and a gigabyte of RAM on the controller, while
Maven may still be running in the same 2\,GB container. The Jenkinsfile
keeps the stages sequential for that reason; on this box, parallel
stages would only compete.
\end{hardware}

\section{Verifying the milestone}

\begin{handson}[the milestone-7 checklist]
\begin{enumerate}
  \item \code{kubectl -n ts get ingress} --- two objects on
  \code{ts.local}; \code{kubectl -n ts get pods -l app=frontend} Ready.
  \item Open \code{http://ts.local/} --- the login page. Then paste
  \code{http://ts.local/courses} directly into the address bar: a page,
  not nginx's 404 --- the SPA fallback works.
  \item Register, log in, create a course: the browser's network tab
  shows \code{/api/...} requests with no \code{OPTIONS} preflights ---
  same origin.
  \item \code{curl -sI http://ts.local/assets/index-*.js | grep
  Cache-Control} --- \code{immutable}; the same for \code{/} --- no
  cache header.
  \item Change any text in a page, push, and watch the pipeline: the
  \code{Frontend image} stage reuses the \code{npm ci} layer, and the
  rollout swaps a 15\,MB image.
\end{enumerate}
\end{handson}

Build~\#9 was the first run of the new stage: \code{docker build}
completed in 90 seconds on the controller, the rollout swapped
\code{frontend:f5d9213} in, and the checklist above passed as written
--- \code{/courses} answered \code{index.html}, \code{/api/courses}
answered the gateway's 401 JSON, the asset carried
\code{immutable} and \code{/} carried no cache header.

\begin{pitfall}[the Ingress that outlived its manifest]
\code{kubectl get ingress} on that day listed \emph{three} objects on
\code{ts.local}: the gateway's, the frontend's --- and
\code{course-service}, eight days old, routing \code{/courses}
straight to the service with no JWT check. Its manifest had been
removed from \code{course-service.yaml} back in Chapter~18 with a
comment saying so, but Chapter~20's rule held: \code{apply} never
deletes what it no longer sees. It took a manual
\code{kubectl -n ts delete ingress course-service}. The lesson is not
the command but the audit: after every milestone, list what the
cluster \emph{has}, not only what git \emph{says} --- until GitOps
pruning (Chapter~27) makes the two the same thing.
\end{pitfall}

\section{Reaching it from anywhere: the tailnet name}

Everything above answered on \code{ts.local} --- a name that exists
only in one laptop's hosts file, on one LAN. ``Live'' means a URL that
resolves on every device you own, and the machine already has one:
Tailscale gives each node a MagicDNS name,
\code{ts-server.tailbfe002.ts.net}, that every peer on the tailnet
resolves without any hosts-file edit. Traefik routes on the
\code{Host} header (Chapter~18), so a name it does not list gets a
404 even though the packets arrive. The fix is one more \code{rules}
entry with the same path list in both Ingress objects
(Appendix~D shows them) --- applied by hand with a filter that
extracts only the Ingress documents from \code{kubectl kustomize}
output, because a full \code{apply -k} would reset every
Deployment's image tag to the manifests' \code{1.0.0-SNAPSHOT} and
roll the whole namespace (Chapter~23's \code{set image} lives
\emph{after} \code{apply} for exactly this reason).

The name then surfaced a bug the LAN had hidden. A registration
through the new host produced its activation mail (Chapter~25's
chain, end to end over the tailnet), but the link inside read
\code{http://localhost:8080/api/auth/activate?code=...}: a host no
phone can reach, and a \code{POST} endpoint that a mail client's
\code{GET} could never have used anyway. \code{EmailService} built
every link --- activate, invite, login --- from string literals. The
correction is the Chapter~16 rule applied to URLs: the public origin
is \emph{configuration}. \code{app.frontend.base-url} defaults to the
Vite dev server in config-server's \code{notification-service.yml};
the manifest overrides it with \code{APP\_FRONTEND\_BASE\_URL} set to
the tailnet name, the same relaxed-binding trick as
\code{SPRING\_MAIL\_HOST}; and the links now target the frontend's
\code{/activate} page, which does the \code{POST}. The two OAuth
variables moved to the same host, which means one more entry under
Google's ``Authorized redirect URIs'' --- the registration that
Chapter~5 called the one thing no manifest can automate.

\begin{pitfall}[a link is an external contract]
Anything that leaves the system --- an email, a webhook payload, a
QR code --- carries a URL that must work from \emph{outside}, so it
can never be a Service name and should never be a literal. The test
is cheap: register from a device that is not the build laptop and
click. On this stack it was the first request through the new
hostname that exposed a bug present since the first commit.
\end{pitfall}

\section{Outside the project}

The pattern --- edge router splits static from API by path, static
comes from a dumb file server --- is universal, with the file server
usually replaced by object storage plus a CDN (S3 + CloudFront, or a
static host such as Netlify) and the split done by the CDN's origin
rules. The build stage stays exactly this Dockerfile's first half.
Two habits transfer directly: hash-and-cache-forever for assets with
an uncached \code{index.html}, and the SPA fallback --- which in
CloudFront terms is ``map 403/404 to \code{/index.html} with status
200'', the same idea with worse ergonomics.

\begin{quiz}
\begin{enumerate}
  \item Why does the frontend image have two \code{FROM} lines, and
  what exactly crosses between them?
  \item A user bookmarks \code{http://ts.local/courses/7} and opens it
  the next morning. Trace the request through Traefik and nginx, and
  say which line makes it succeed.
  \item Both Ingress objects declare \code{host: ts.local}. How does
  Traefik decide that \code{/api/courses} goes to the gateway and
  \code{/assets/x.js} to nginx?
  \item Explain why the gateway's \code{CorsConfig} became irrelevant
  in the cluster while remaining necessary on a laptop.
  \item What does \code{--group-add "\$(stat -c \%g
  /var/run/docker.sock)"} fix, and why is the value computed instead
  of written down?
  \item The activation mail's link was wrong for months and nobody
  noticed. Name the two separate defects in
  \code{http://localhost:8080/api/auth/activate?code=...}, and say
  which environment (laptop, compose, cluster) each one would have
  broken.
\end{enumerate}
\end{quiz}
